% Part: turing-machines
% Chapter: machines-computations
% Section: unary-numbers

\documentclass[../../../include/open-logic-section]{subfiles}

\begin{document}

\olfileid{tur}{mac}{una}
\olsection{نمایش یکانی اعداد}

\begin{explain}
ماشین‌های تورینگ روی دنباله‌های نمادهای نوشته‌شده بر نوارشان کار
می‌کنند. بسته به الفبای ماشین تورینگ، این دنباله‌های نمادها می‌توانند
نمایندهٔ ورودی‌ها و خروجی‌های گوناگون باشند. البته ماشین‌های تورینگی
که توابع \emph{حسابی}، یعنی توابع اعداد طبیعی، را محاسبه می‌کنند
اهمیت ویژه‌ای دارند. راهی ساده برای نمایش اعداد صحیح مثبت، کدکردن
آن‌ها به‌صورت دنباله‌هایی از نماد یگانهٔ~$\TMstroke$ است. اگر
$n\in\Nat$، بگذارید $\TMstroke^n$ برای $n=0$ دنبالهٔ تهی و در غیر
این صورت دنباله‌ای متشکل از دقیقاً $n$ عدد $\TMstroke$ باشد.
\end{explain}

\begin{defn}[محاسبه]
ماشین تورینگ~$M$ تابع $f\colon\Nat^k\to\Nat$ را \emph{محاسبه
می‌کند} اگر و تنها اگر برای همهٔ $n_1$، \dots،~$n_k\in\Nat$، $M$ روی ورودی
\[
\TMstroke^{n_1} \TMblank \TMstroke^{n_2} \TMblank \dots \TMblank \TMstroke^{n_k}
\]
با خروجی $\TMstroke^{f(n_1,\dots,n_k)}$ متوقف شود.
\end{defn}

\begin{prob}
  تعریف کنید ماشین تورینگ~$M$ چه هنگام تابع
  $f\colon\Nat^k\to\Nat^m$ را محاسبه می‌کند.
\end{prob}

\begin{ex}\ollabel{ex:adder}
\emph{جمع:}
ماشینی بسازیم که تابع $f(n,m)=n+m$ را محاسبه کند. به ماشینی نیاز داریم
که با دو بلوک $\TMstroke$ به طول‌های $n$ و~$m$ روی نوار آغاز کند و
با یک بلوک متشکل از $n+m$ عدد $\TMstroke$ متوقف شود. دو بلوک ورودی
$\TMstroke$ با یک~$\TMblank$ جدا می‌شوند؛ پس یک روش آن است که در
خانهٔ دارای $\TMblank$ خطی بنویسیم و آخرین~$\TMstroke$ را پاک کنیم.
\begin{figure}\[
\begin{tikzpicture}[->,>=stealth',shorten >=1pt,auto,node distance=2.8cm,
                    semithick]
  \tikzstyle{every state}=[fill=none,draw=black,text=black]

  \node[initial,state]         (A)              {$q_0$};
  \node[state]         (B) [right of=A] {$q_1$};
  \node[state]         (C) [right of=B] {$q_2$};

  \path (A) edge node {\TMtrans{\TMblank}{\TMstroke}{\TMstay}} (B)
           edge [loop above] node {\TMtrans{\TMstroke}{\TMstroke}{\TMright}} (A)
        (B) edge [loop above] node {\TMtrans{\TMstroke}{\TMstroke}{\TMright}} (B)
            edge node {\TMtrans{\TMblank}{\TMblank}{\TMleft}} (C)
        (C) edge [loop above] node {\TMtrans{\TMstroke}{\TMblank}{\TMstay}} (C);
\end{tikzpicture}
\]\caption{ماشینی که $f(x,y)=x+y$ را محاسبه می‌کند}
\ollabel{fig:adder}
\end{figure}
\end{ex}

\begin{prob}
پیکربندی‌های ماشین \olref[tur][mac][una]{ex:adder} را برای ورودی
$\tuple{3,2}$ گام‌به‌گام دنبال کنید. اگر ماشین $0+0$ را محاسبه کند
چه رخ می‌دهد؟
\end{prob}

\begin{explain}
در \olref[rep]{ex:doubler} نمونه‌ای از ماشین تورینگ ارائه کردیم که
دنباله‌ای از~$\TMstroke$ها را ورودی می‌گیرد و با دنباله‌ای دو برابر
بلندتر از~$\TMstroke$ها روی نوار متوقف می‌شود---ماشین دوبرابرکن.
اما چون خروجی در چپ بلوک دوبرابرشدهٔ~$\TMstroke$ها تعدادی
$\TMblank$ دارد، برخلاف آنچه شاید تصور کرده باشید در واقع تابع
$f(x)=2x$ را محاسبه نمی‌کند. دو راه اصلاح آن را شرح می‌دهیم.
\end{explain}

\begin{ex}
ماشین \olref{fig:doubler-disc} تابع $f(x)=2x$ را محاسبه می‌کند. این
ماشین به‌جای آنکه مانند ماشین \olref[rep]{ex:doubler} به ازای هر
$\TMstroke$ ورودی، ورودی را پاک کند و دو $\TMstroke$ در انتهای راست
بنویسد، به ازای هر $\TMstroke$ ورودی یک~$\TMstroke$ به راست می‌افزاید.
باید محل پایان ورودی را پیگیری کند؛ پس یک~$\TMblank$ میان ورودی و
خط‌های افزوده باقی می‌گذارد که در پایان با~$\TMstroke$ پر می‌کند.
همچنین باید جای خود را در ورودی «به یاد بسپاریم»، پس موقتاً یک
$\TMstroke$ از بلوک ورودی را با~$\TMblank$ جایگزین می‌کنیم.
  \begin{figure}
\[
\begin{tikzpicture}[->,>=stealth',shorten >=1pt,auto,node distance=2.8cm,
                    semithick]
  \tikzstyle{every state}=[fill=none,draw=black,text=black]

  \node[initial,state] (0)              {$q_0$};
  \node[state]         (1) [above of=0] {$q_1$};
  \node[state]         (2) [above of=1] {$q_2$};
  \node[state]         (3) [right of=2] {$q_3$};
  \node[state]         (4) [below of=3] {$q_4$};
  \node[state]         (5) [below of=4] {$q_5$};
  \node[state]         (6) [above right of=4] {$q_6$};
  \node[state]         (7) [below of=6] {$q_7$};
  \node[state]         (8) [below of=7] {$q_8$};

  \path (0) edge node {\TMtrans{\TMstroke}{\TMblank}{\TMright}} (1)
%    (0) edge node {\TMtrans{\TMblank}{\TMblank}{\TMstay}} (h)
    (1) edge [loop left] node {\TMtrans{\TMstroke}{\TMstroke}{\TMright}} (1)
      edge node {\TMtrans{\TMblank}{\TMblank}{\TMright}} (2)
    (2) edge [loop above] node {\TMtrans{\TMstroke}{\TMstroke}{\TMright}} (2)
        edge node {\TMtrans{\TMblank}{\TMstroke}{\TMleft}} (3)
    (3) edge [loop above] node {\TMtrans{\TMstroke}{\TMstroke}{\TMleft}} (3)
        edge node[left] {\TMtrans{\TMblank}{\TMblank}{\TMleft}} (4)
    (4) edge        node[left] {\TMtrans{\TMstroke}{\TMstroke}{\TMleft}} (5)
    (5) edge [loop below] node {\TMtrans{\TMstroke}{\TMstroke}{\TMleft}} (5)
        edge              node {\TMtrans{\TMblank}{\TMstroke}{\TMright}} (0)
    (4) edge      node[sloped] {\TMtrans{\TMblank}{\TMstroke}{\TMright}} (6)
    (6) edge              node {\TMtrans{\TMblank}{\TMstroke}{\TMright}} (7)
    (7) edge [loop right] node {\TMtrans{\TMstroke}{\TMstroke}{\TMright}} (7)
        edge              node {\TMtrans{\TMblank}{\TMblank}{\TMleft}} (8)
    (8) edge [loop right] node {\TMtrans{\TMstroke}{\TMblank}{\TMstay}} (8);
    \end{tikzpicture}
\]
\caption{ماشینی که $f(x)=2x$ را محاسبه می‌کند}
\ollabel{fig:doubler-disc}
\end{figure}
\end{ex}

\begin{ex}\ollabel{ex:mover}
راه دوم برای محاسبهٔ $f(x)=2x$ این است که ماشین دوبرابرکن اصلی را
نگه داریم، اما در پایان حالت‌ها و دستورهایی بیفزاییم که بلوک
دوبرابرشدهٔ خط‌ها را به انتهای چپ نوار منتقل کنند. ماشین
\olref{fig:mover} دقیقاً این بخش آخر را انجام می‌دهد: روی نواری متشکل
از بلوکی از~$\TMblank$ها و سپس بلوکی از~$\TMstroke$ها آغاز می‌کند
(سر در هر جایی از بلوک~$\TMblank$ها قرار دارد)، $\TMstroke$ها را
یکی‌یکی پاک می‌کند و آن‌ها را در آغاز نوار می‌نویسد. برای تشخیص زمان
پایان کار، نخست انتهای بلوک $\TMstroke$ها را با نماد~$\TMendtape$
علامت می‌زند که در پایان پاک می‌شود. شماره‌گذاری حالت‌ها را از~$q_6$
آغاز کرده‌ایم تا بتوان آن‌ها را به ماشین دوبرابرکن افزود. تنها به
دستور اضافی
$\delta(q_0,\TMblank)=\tuple{q_6,\TMblank,\TMstay}$ نیاز دارید؛ یعنی
پیکانی از~$q_0$ به~$q_6$ با برچسب
$\TMtrans{\TMblank}{\TMblank}{\TMstay}$. (یک مشکل ظریف وجود دارد:
ماشین حاصل برای ورودی~$x=0$ کار نمی‌کند. آن را به‌عنوان تمرین
می‌گذاریم.)
\begin{figure}
  \[
  \begin{tikzpicture}[->,>=stealth',shorten >=1pt,auto,node distance=2.8cm,
                      semithick]
    \tikzstyle{every state}=[fill=none,draw=black,text=black]
    \node[initial,state] (6)              {$q_6$};
    \node[state]         (7) [right of=6] {$q_7$};
    \node[state]         (8) [right of=7] {$q_8$};
    \node[state]         (9) [below of=8] {$q_9$};
    \node[state]         (10) [left of=9] {$q_{10}$};
    \node[state]         (11) [left of=10]  {$q_{11}$};
    \node[state]         (12) [below of=11]  {$q_{12}$};
    \node[state]         (13) [right of=12] {$q_{13}$};
    \node[state]         (14) [right of=13] {$q_{14}$};
    \path
    (6)  edge node {\TMtrans{\TMblank}{\TMblank}{\TMright}} (7)
    (7)  edge [loop above] node {\TMtrans{\TMblank}{\TMblank}{\TMright}} (7)
         edge node {\TMtrans{\TMstroke}{\TMstroke}{\TMright}} (8)
    (8)  edge [loop above] node {\TMtrans{\TMstroke}{\TMstroke}{\TMright}} (8)
         edge node {\TMtrans{\TMblank}{\TMendtape}{\TMleft}} (9)
    (9)  edge [loop right] node {\TMtrans{\TMstroke}{\TMstroke}{\TMleft}} (9)
         edge node {\TMtrans{\TMblank}{\TMblank}{\TMright}} (10)
    (10) edge node[above] {\TMtrans{\TMstroke}{\TMblank}{\TMleft}} (11)
    (11) edge [loop above] node {\TMtrans{\TMblank}{\TMblank}{\TMleft}} (11)
         edge node[left] {\begin{tabular}{@{}l@{}}
          \TMtrans{\TMendtape}{\TMendtape}{\TMright}\\
          \TMtrans{\TMstroke}{\TMstroke}{\TMright}
         \end{tabular}} (12)
    (12) edge node[] {\TMtrans{\TMblank}{\TMstroke}{\TMright}} (13)
    (13) edge [loop below] node {\TMtrans{\TMblank}{\TMblank}{\TMright}} (13)
         edge node[sloped] {\TMtrans{\TMstroke}{\TMblank}{\TMleft}} (11)
    (13) edge node {\TMtrans{\TMendtape}{\TMblank}{\TMstay}} (14);
  \end{tikzpicture}
  \]
  \caption{انتقال بلوکی از $\TMstroke$ها به چپ}
  \ollabel{fig:mover}
\end{figure}
\end{ex}

\begin{prob}
در \olref[tur][mac][una]{ex:mover} ماشینی مرکب از ماشین دوبرابرکنِ
\olref[tur][mac][una]{fig:doubler-disc} و ماشین جابه‌جاکنِ
\olref[tur][mac][una]{fig:mover} توصیف کردیم. اگر این ماشین ترکیبی را
روی ورودی~$x=0$، یعنی نواری تهی، آغاز کنید چه رخ می‌دهد؟ چگونه ماشین
را اصلاح می‌کنید تا در این حالت با خروجی~$2x=0$ متوقف شود؟ (باید
بتوانید با افزودن یک حالت و یک گذار چنین کنید.)
\end{prob}

\begin{prob}
\emph{تفریق:} ماشین تورینگی طراحی کنید که با گرفتن دو رشتهٔ ناتهی از
خط‌ها به طول $n$ و~$m$، که $n>m$، تابع $f(n,m)=n-m$ را محاسبه کند.
\end{prob}

\begin{prob}
\emph{برابری:} ماشین تورینگی برای محاسبهٔ تابع زیر طراحی کنید:
\[
\fn{equality}(n,m) = 
\begin{cases}
  \text{1} & \text{اگر $n = m$} \\
  \text{0} & \text{اگر $n \neq m$}
\end{cases}
\]
که در آن $n$ و~$m\in\PosInt$.
\end{prob}

\begin{prob}
ماشین تورینگی برای محاسبهٔ تابع $\min(x,y)$ طراحی کنید که در آن $x$
و~$y$ اعداد صحیح مثبت‌اند و روی نوار با رشته‌هایی از $\TMstroke$ که
یک $\TMblank$ آن‌ها را جدا می‌کند نمایش داده شده‌اند. می‌توانید در
الفبای ماشین از نمادهای اضافی استفاده کنید.

تابع $\min$ کوچک‌ترین مقدار را از میان آرگومان‌هایش برمی‌گزیند؛ پس
$\min(3,5)=3$، $\min(20,16)=16$ و $\min(4,4)=4$، و به همین ترتیب.
\end{prob}

\begin{defn}
  ماشین تورینگ~$M$ تابع جزئیِ $f\colon\Nat^k\pto\Nat$ را محاسبه
  می‌کند اگر و تنها اگر
  \begin{enumerate}
    \item اگر $f(n_1,\dots,n_k)=m$، آنگاه $M$ روی ورودی
    $\TMstroke^{n_1}\concat\TMblank\concat
    \dots\concat\TMblank\concat\TMstroke^{n_k}$ با خروجی
    $\TMstroke^{m}$ متوقف شود.
    \item اگر $f(n_1,\dots,n_k)$ تعریف‌نشده باشد، آنگاه $M$ اصلاً
    متوقف نشود یا با خروجی‌ای متوقف شود که یک بلوک یگانه از
    $\TMstroke$ها نیست.
  \end{enumerate}
\end{defn}

\end{document}
